% !PRINTED 691
\weberpage{691}
%<<CONT
und $2\nu - n$ reelle vorkommen. Demnach sind von den $n$ Functionen $y$ [§. 189, (2)]

\weq{2}{y_k = \alpha_{1,k}x_1 + \alpha_{2,k}x_2 + \cdots + \alpha_{n,k}x_n.}

$2\nu - n$ reell, $2n-2\nu$ paarweise conjugirt. Wir wollen unter den $n$ Variablen $u$ reelle lineare Verbindungen der $x$ verstehen, nämlich die $2\nu-n$ reellen $y$ und die $2n-2\nu$ Bestandtheile der conjugirt imaginären Paare der $y$, so dass, wenn $y, y'$ ein imaginäres Paar ist,
\[
2u_1 = y+y', \quad 2iu_2 = y-y'
\]
oder

\weq{3}{y = u_1+iu_2, \quad y' = u_1-iu_2}

gesetzt wird. Wenn man in der Functionaldeterminante der $y$ auf jedes Paar conjugirt imaginärer Zeilen zweimal den elementaren Determinantensatz von der Addition der Zeilen anwendet (Bd. I, §. 22), so ergiebt sich

\[
\begin{aligned}
\Sigma \pm \frac{\partial u_1}{\partial x_1}\frac{\partial u_2}{\partial x_2}\cdots\frac{\partial u_n}{\partial x_n}
&= \frac{1}{(-2i)^{n-\nu}}\Sigma\pm\frac{\partial y_1}{\partial x_1}\frac{\partial y_2}{\partial x_2}\cdots\frac{\partial y_n}{\partial x_n} \\
&= \frac{A}{(-2i)^{n-\nu}},
\end{aligned}
\]

worin $A$ die in §. 189, (3) festgesetzte Bedeutung hat.

Die Körperdiscriminante ist positiv oder negativ, je nachdem die Anzahl der conjugirt imaginären Paare, also $n-\nu$, gerade oder ungerade ist, und demnach erhält man, wenn mit $\pm\varDelta$ der absolute Werth der Körperdiscriminante bezeichnet wird, nach §. 189, (4), (16):

\weq{4}{\mathrm{Lim}\,\frac{T}{t} = \frac{2^{n-\nu}}{w\,N(\mathfrak{a})\sqrt{\pm\varDelta}}\int\int\cdots\int du_1\,du_2\ldots du_n.}

Es ist also weiter noch das Integral

\[
U = \int\int\cdots\int du_1\,du_2\ldots du_n
\]

zu behandeln.

Aus den Gleichungen §. 189, (10) erhält man zu jedem den Bedingungen (15) genügenden Werthsysteme der Variablen $\xi_1, \xi_2, \ldots, \xi_\nu$ ein bestimmtes endliches System der Werthe für $z_1, z_2, \ldots, z_\nu$, und daraus nach §. 189, (7) für jedes $y$ einen von Null verschiedenen absoluten Werth. Setzt man im Falle eines imaginären Paares (3)

\weq{5}{u_1 = e^{1/2\,z}\cos\vartheta, \quad u_2 = e^{1/2\,z}\sin\vartheta,}

\weq{6}{y_1y_1' = u_1^2+u_2^2 = e^z,}
%CONT>>
